\documentclass[a4paper,twoside,12pt]{book} %aaaaa čo je to vlastne za prácu ta bakalárka? book/report?
% užitočná linka http://en.wikibooks.org/wiki/LaTeX/
% http://www.scribd.com/doc/2633955/Writing-a-thesis-with-LaTeX
% http://www.eng.cam.ac.uk/help/tpl/textprocessing/reports.html
% documentclass možno volať aj s [openright] ale asi nefunguje
%%načítanie balíkov

%lintex!

% zlý, latex, zlý! :)
%\usepackage[utf8x]{inputenc} %používať utf8x je ževraj nebezpečné!!
%\usepackage[T1]{fontenc}

%%xelatex!
\usepackage[xetex]{graphicx}
\usepackage{fontspec,xunicode}
\usepackage{amsmath}
%\usepackage{unicode-math}
\defaultfontfeatures{Mapping=tex-text,Scale=MatchLowercase,Ligatures={Common}}
%\setmainfont{Nimbus Roman No9 L}
%\setmonofont{LMMonoLt10}
\setmainfont{Junicode}
%\setmainfont{Gentium Basic}
%\setmainfont{Asana-Math}
\newICUfeature{LigType}{disc}{+dlig}
%\usepackage[bibstyle=numeric, citestyle=numeric-comp, sorting=none, backend=biber]{biblatex}
%\bibliography{literatura}

%http://www.charlietanksley.net/philtex/the-unicode-math-package-for-xelatex-and-the-stix-fonts/
%\setsansfont[Mapping=tex-text]{Arial}
%\setmonofont{Courier New}
%\setmathfont{Asana-Math}

%\setmathfont{STIXGeneral}

% \usepackage{microtype}

\usepackage[english,czech,slovak]{babel}
% \usepackage{slovak}
%%slovenská deľba slov
%http://ctan.org/pkg/polyglossia
% \usepackage{polyglossia}
% \setdefaultlanguage{slovak}
% \usepackage[slovak]{babel}

% \usepackage[numbers]{natbib}
% \usepackage[round]{natbib}

\usepackage[german=quotes,style=german]{csquotes}
% \usepackage{csquotes}
% \DeclareQuoteStyle{czech}%     [variant]{style}
%   {\quotedblbase }%                   opening outer mark
%   {\textquotedblleft}%                  closing outer mark
%   {\textquoteleft}%                      opening inner mark
%   {\textquoteright}%    
% %\DeclarePlainStyle{,,}{''}{,}{'}
% \setquotestyle{czech}
% \usepackage[style=authoryear-comp,bibstyle=authoryear,backend=biber]{biblatex}
\usepackage[
   backend=biber      % if we want unicode 
  ,style=authoryear-comp % iso-authoryear or iso-numeric for numeric citation method
 % ,bibstyle=iso          
  ,babel=other        % to support multiple languages in bibliography
  %,sortlocale=cs_CZ   % locale of main language, it is for sorting
  ,bibencoding=UTF8   % this is necessary only if bibliography file is in different encoding than main document
]{biblatex}
% \input{bibiso}
% \input{bibloc}
% style
% authoryear-comp
% verbose-trad1
% bibstyle
% authoryear
\def\citet#1{\textcite{#1}}
\def\citep#1{\parencite{#1}}
\addbibresource{literatura.bib}
 
\usepackage{graphics}
\usepackage{graphicx}
\usepackage{subfig}
\usepackage{array}

%%nastavenie okrajov
\usepackage[a4paper]{geometry}
%toto je fakt nechutné odsádzanie..
%\geometry{top=2.5cm,bottom=2.5cm,left=3.5cm,right=2cm}
\geometry{top=2.5cm,bottom=2.5cm,left=3.5cm,right=2.0cm}
%%riadkovanie primárneho textu
\usepackage{setspace}
%\singlespacing
\onehalfspacing
%\doublespacing
%\setstretch{1.1}

\usepackage{changepage}


%%package, ktoré sú užitočné, ale zatiaľ som ich nepoužil
%\usepackage{lastpage}
%\usepackage{layouts}
%\usepackage[]{fp}
%\usepackage{listings-ext}
%\lstset{language=C++}
%\usepackage[comma,authoryear]{natbib}
%

  
\usepackage{url}
\usepackage[a4paper=true,hypertex]{hyperref}

%%definície konštánt
\def \nazovUniverzity{Univerzita Komenského v Bratislave}
\def \NazovUniverzity{\textbf{\Large \sc \nazovUniverzity}} %\sc pre smallcaps
\def \nazovFakulty{Fakulta matematiky, fyziky a informatiky}
\def \NazovFakulty{\textbf{\large \sc \nazovFakulty}}
\def \nazovDiela{Vizualizácia algoritmov a dátových štruktúr}
\def \NazovDiela{\LARGE{\em \nazovDiela}}
\def \podnazovDiela{nie je}
\def \PodazovDiela{\large{\em \podnazovDiela}}
\def \typPrace{Bakalárska práca}
\def \TypPrace{\large{\typPrace}}
\def \autor{Viktor Tomkovič}
\def \veduci{Mgr. Jakub Kováč}
\def \rok{2011}
\def \miestoRok{Bratislava, \rok}
\def \cisloOdboru{2508}
\def \katedra{Katedra Informatiky FMFI}
\def \odbor{Informatika}
\def \program{Informatika}

%%škaredé veci
%\def \lavyokraj{\oddsidemargin}%\printinunitsof{cm}\prntlen{\textwidth}}%\prntlen{\oddsidemargin}}
%\def \pravyokraj{\the\pdfpagewidth-\the\oddsidemargin+\the\textwidth}
%\FPadd{}{\the\pdfpagewidth}{\the\textwidth}

%%definície príkazov
%citat: argumenty sú citát a meno autora
%\citat{Live as if you were to die tomorrow\\Learn as if you were to live forever.}{Mahátma Gándhí}
% Copy from one, it’s plagiarism; copy from two, it’s research. Wilson Mizner
\newcommand{\citat}[2]{\begin{flushright}\textsl{#1}\\#2\end{flushright}}
%strednp je vystredenie textu na nepárnej strane bez ohľadu na okraje(!!!), parametrom je text
\newcommand{\stredp}[1]{\begin{adjustwidth}{+0.50cm}{+0.00cm} \begin{center} #1 \end{center} \end{adjustwidth}}
%stredp je obdoba predošlého príkazu len pre párne strany
\newcommand{\strednp}[1]{\begin{adjustwidth}{+0.0cm}{+0.50cm} \begin{center} #1 \end{center} \end{adjustwidth}}

%To supress only the word chapter, add this in the preamble:
%\renewcommand{\chaptername}{} %nefunguje
%If you also want to supress the numbering then add
%\renewcommand{\thechapter}{}

%\usepackage{epsfig}
%\usepackage{url}

% \DeclareTextCommandDefault\clqq{\leavevmode
%   \set@low@box{„}%
%   \setbox1=\hbox{l\/}\dimen1=\wd1
%   \setbox1=\hbox{l}\advance\dimen1 by -\wd1
%   \ifdim\dimen1>0pt\kern-.1em\box0\kern.1em
%     \else\kern.1em\box0\kern-.1em\fi\nobreak\hskip0pt}
% \DeclareTextCommandDefault\crqq{\edef\@SF{\spacefactor\the\spacefactor}%
%   \nobreak\kern-.07em\hbox{“}\kern.07em\@SF\relax}

\def\uv#1{„#1“}
%  \newcommand{\uv}[1]{\clqq #1\crqq}
% \def\uv#1{\clqq #1\clrqq}
\def\todo#1{\par\smallskip{\scriptsize {\bf TODO: }#1}\par\smallskip}

\def\ins{\mathop{\mathit{insert}}}
\def\find{\mathop{\mathit{find}}}
\def\delete{\mathop{\mathit{delete}}}
\def\dec{\mathop{\mathit{decreaseKey}}}
\def\delmin{\mathop{\mathit{deleteMin}}}
\def\meld{\mathop{\mathit{meld}}}
\def\rank{\mathop{\mathit{rank}}}
\def\L{\mathop{\mathit{left}}}
\def\R{\mathop{\mathit{right}}}
\def\reverse{\mathop{\mathit{reverse}}}
\def\Bp{$\hbox{\rm B}\!^+\!$}

\def\null{\texttt{NULL}}
\def\makeset{$\mathop{\mathit{makeset}}(x)$}
\def\find{$\mathop{\mathit{find}}(x)$}
\def\union{$\mathop{\mathit{union}}(x, y)$}

\def\k{w}
%\def\kluc{kľúč}
\def\put{$\mathop{\mathit{insert}}(\k)$}
\def\find{$\mathop{\mathit{find}}(\k)$}
\def\delete{$\mathop{\mathit{delete}}(\k)$}
\def\trie{trie}
\def\uz{\hbox{\tt\$}}


\catcode`\"=\active \def "{\begingroup„\def "{\endgroup“}}
\begin{document}
% \catcode`\"=\active \def "{\begingroup\clqq\def "{\endgroup\crqq}}

\frontmatter

\thispagestyle{empty}
\noindent
\strednp{
\NazovUniverzity\\ 
\NazovFakulty
}
\vfill
\strednp{
\NazovDiela
%% \centerline{\PodazovDiela}
\mbox{}\\
\bigskip
\TypPrace
}
\vfill
\strednp{\rok \hfill{\bf \autor}}
\newpage

\thispagestyle{empty}
\noindent
\strednp{\NazovUniverzity\\ \NazovFakulty}
\vfill
\strednp{\NazovDiela
%% \centerline{\PodazovDiela}
\mbox{}\\
\bigskip
\TypPrace
}
\vfill
\begin{tabular}{ l l }
\textbf{Študijný program:} & \program\\
\textbf{Študijný odbor:} & \cisloOdboru\ \odbor\\
\textbf{Školiace pracovisko:} & \katedra\\
\textbf{Vedúci práce:} &  \veduci
\end{tabular}
\bigskip\\
\bigskip\\
\bigskip\\
\bigskip\\
\strednp{\miestoRok \hfill{\bf \autor}}
\newpage

\noindent
~\vfill

\section*{Poďakovanie}
Veľká vďaka za to, že mi pomohli patrí môjmu školiteľovi Kubovi Kováčovi a 
mojim rodičom.\\
\bigskip\\
\newpage

\chapter*{Abstrakt}
Táto práca sa zaoberá vizualizáciou dátových štruktúr. Nadväzuje na bakalársku 
prácu Jakuba Kováča (2007) a rozvíja kompiláciu o ďalšie dátové štruktúry. V 
prvej časti sme tieto dátové štruktúry, konkrétne union-find, lexikografický 
strom a sufixový strom, definovali a konkrétne algoritmy popísali. V druhej 
časti sme spravili návrh aplikácie a popísali implementáciu. Práca je 
prehľadom dátových štruktúr s priloženým Java appletom na CD ako pomôckou pri 
vizualizácií.\\
Kľúčové slová: vizualizácia, ADŠ, algoritmy a dátové štruktúry, union-find, 
stringológia, lexikografický strom, trie, sufixový strom.

\newpage

\chapter*{Abstract}
This work is about visualization of algorithms and data structures. Continues 
the bachelor thesis of Jakub Kováč (2007) and extends the compilation by 
another data structures. In the first part are the structures, namely 
union-find, lexicographical tree (trie) and suffix tree, defined and the 
algorithms described. The second part is a description and a implementation of 
the software. The work is an overview of data structures with a software CD 
included.\\
Keywords: visualization, ADS, algorithms and data structures, union-find, 
stringology, loxicographical tree, trie, suffix tree.
\newpage

\mbox{}
\newpage

\tableofcontents\newpage
%\listoffigures
%\listoftables


\mainmatter

%\part{Vizualizácia a dátové štruktúry}

\input{uvod}

% \indent Dátové štruktúry a algoritmy tvoria základnú, prvotnú časť výučby informatiky.
% Vizualizácia algoritmov a dátových štruktúr (ďalej len ako "vizualizácia") je grafické znázornenie, 
% ktoré abstrahuje spôsob ako algoritmus a dátové štruktúry pracujú od ich vnútornej reprezentácie a umiestení v pamäti.
% Je teda vyhľadávaná a všeobecne rozšírená pomôcka pri výučbe.
% Výsledky výskumov ohľadne jej efektívnosti sa rôzňa, od stavu "nezaznamenali sme výrazné zlepšenie" po "je viditeľné zlepšenie".
% \cite{Shaffer:2010:AVS:1821996.1821997}
% Zaujímavé je pozorovanie, že určovanie si vlastného tempa pri vizualizácií je veľká pomôcka, 
% oproti tomu ukazovaniu pseudokódu alebo nemožnosť určovania tempa, napríklad animácia bez možnosti pozastavenia,
% takmer žiadne zlepšenie neprináša.
% \cite{Shaffer:2010:AVS:1821996.1821997}\cite{Saraiya:2004:EFA:971300.971432}
% Rozmach vizualizačných algoritmov priniesla najmä Java a jej fungovanie, bez viazanosti na konkrétny operačný systém. 
% Kvalita vizualizácií sa líši a keďže ide o ľahko naprogramovateľné programy, je ich veľa a sú pomerne nekvalitné. 
% V takomto množstve je ťažké nájsť kvalitné vizualizácie. 
% Zbieraním a analyzovaním kvality sa venuje skupina AlgoViz, ktorá už veľa rokov funguje na portále algoviz.org.
% 
% Táto práca opisuje algoritmy na dátových štruktúrach: dynamické množiny s operáciami "nájdi"\ a "spoj" 
% (v texte označované union-find), lexikografický strom a suffix tree.
% Union-find je stále skúmaný problém.\cite{Wu:2005}\cite{Aho:1974:DAC:578775}\cite{Yao:1985:EPP:3674.3682}\cite{wiki:uf} 
% V dôkaze Tarjana o najhoršom prípade stále sa zatiaľ chyba nenašla.
% Priemerný prípad je zaujímavejší, práca z roku 2005 ukazuje, že union--find problém je lineárny.\cite{Wu:2005}
% 
% O stringológií som ešte nečítal, takže súčasný stav neviem.
% 

% Uvedené dátové štruktúry sú rozšírené a preto považujem za vhodné spraviť ich kvalitnú vizualizáciu a kompiláciu. 
% Zdôrazňujem kvalitnú, lebo na spomínanom portále algoviz.org, v čase písania žiadna dobrá vizualizácia union-find problému nie je.
% Motiváciou je najmä s ostatnými spoluautormi Kubom Kováčom, Katkou Kotrlovou, Tatianou Tóthovou a Paľom Lukčom vybudovať 
% vizualizačný software, ktorý poslúži ako dobrá pomôcka učiteľom a študentom informatiky nielen na Slovensku, ale aj v zahraničí.
% (ako mám napísať, že Katreniaková, Kováč už niečo podobné na našej fakulte robili a Kotrlová momentálne robí?)


\chapter{Union-find problém}\label{chap:uf}

V niektorých aplikáciach potrebujeme udržiavať prvky rozdelené do skupín
(disjunktných množín), pričom skupiny sa môžu zlučovať a my potrebujeme
pre daný prvok efektívne zistiť, do ktorej skupiny patrí. Predpokladáme,
že každá množina $S$ je jednoznačne určená jedným svojim zástupcom $x\in S$
a potrebujeme implementovať nasledovné tri operácie:

\begin{itemize}
\item $\makeset(x)$ -- vytvorí novú množinu $S=\{x\}$ s jedným prvkom; %,
                   %ktorý nepatrí do žiadnej inej množiny;
\item $\union(x, y)$ -- ak $x, y$ sú zástupcovia množín $S$ a $T$,
                 $\union$ vytvorí novú množinu $S\cup T$,
                 pričom $S$ aj $T$ zmaže. Zástupcom novej množiny $S\cup T$
                 je $x$ alebo $y$.
\item $\find(x)$ -- nájde zástupcu množiny, v ktorej sa 
                prvok $x$ nachádza.
\end{itemize}
 
\section{Použitie}\label{sec:uf-pouzitie}

Union-find sa dá použiť na reprezentáciu neorientovaného grafu,
do ktorého pridávame hrany a odpovedáme na otázku: "Sú dané dva
vrcholy spojené nejakou cestou?" (t.~j.\ sú v rovnakom komponente súvislosti?)
Medzi najznámejšie aplikácie patria Kruskalov algoritmus na nájdenie najlacnejšej
kostry \citep{kruskal} a unifikácia \citep{unif}.

\citet{cholesky} ukázali, ako sa dá union-find použiť pri Choleského dekompozícií
riedkych matíc. Autori navrhli efektívny algoritmus, ktorý zistí počet nenulových
prvkov v každom riadku a stĺpci výslednej matice, čo slúži na efektívnu alokáciu
pamäte.

Pre offline verziu úlohy, kde sú všetky operácie dopredu známe, \citet{offline-uf}
navrhli lineárny algoritmus. Článok obsahuje tiež viacero aplikácií v teoretickej
informatike.

\section{Popis}
Dátová štruktúra \emph{union-find} sa reprezentuje ako les, kde každý strom zodpovedá
jednej množine a korene stromov sú zástupcovia množín. Pri implementácií si
stačí pre každý prvok $x$ udržiavať smerník $p(x)$ na jeho otca
(pre koreň je $p(x)=\nul$).

\paragraph{Vytváranie prvkov.}%
Operácia $\makeset(x)$ teda vytvorí nový prvok $x$ a nastaví $p(x) = \nul$. 

\paragraph{Hľadanie zástupcu.}%
Operáciu $\find(x)$ vykonáme tak, že budeme sledovať cestu po smerníkoch, až 
kým nenájdeme zástupcu. 

\paragraph{Spájanie množín.}%
Operáciu $\union(x, y)$ ide najjednoduchšie vykonať tak, že presmerujeme smerník 
$p(y)$ na prvok $x$, teda $p(y) \gets x$. 
Môžeme ľahko pozorovať, že takýto \emph{naivný} spôsob je neefektívny, 
lebo nám operácia $\find(x)$ v najhoršom prípade, na $n$ prvkoch, trvá $\Omega(n)$ 
krokov. 

\bigskip
Existujú dva prístupy ako zlepšiť operácie a tým aj zrýchliť ich vykonanie. 
Sú to: heuristika \emph{spájanie podľa ranku} a rôzne heuristiky na 
\emph{kompresiu cesty}. 

\section{Heuristika na spájanie}\label{sec:komp-union}

Prvá heuristika pre každý vrchol $x$ udržuje hodnotu $\rank(x)$,
ktorá určuje najväčšiu možnú hĺbku podstromu s koreňom $x$.
Pri o\-pe\-rá\-cií $\makeset(x)$ zadefinujeme $\rank(x) = 1$. 
Pri o\-pe\-rá\-cií $\union(x, y)$ porovnáme $\rank(x)$ a $\rank(y)$
a vždy napojíme strom s menším rankom pod strom s väčším rankom.
Ak majú oba stromy rovnaký rank, napojíme povedzme $x$ pod $y$
a $\rank(y)$ zvýšime o 1.
%aby sme zistili, ktorý zástupca predstavuje menší strom. 
% Smerník tohto zástupcu potom napojíme 
% na zástupcu s výšším rankom. Zástupca novej množiny bude ten s vyšším rankom. 
% Ak sú oba ranky rovnaké, vyberieme ľubovoľného zo zástupcov $x$ a $y$, 
% jeho rank zvýšime o jeden a smerník ostatného zástupcu bude ukazovať 
% na tohto zástupcu. %Zástupcom novej množiny bude vybratý zástupca. 

\section{Heuristiky na kompresiu cesty}\label{sec:komp-cesta}
      
\begin{figure}
\centering
\subfloat[Pred vykonaním kompresie.]{
\label{img:komp1}
% ,width=0.4\textwidth
\includegraphics[height=0.24\textheight]{obrazky/komp1.png}
}
\qquad
\subfloat[Úplná kompresia.]{
\label{img:komp2}
\includegraphics[height=0.12\textheight]{obrazky/komp2.png}
}

\subfloat[Delenie cesty.]{
\label{img:komp3}
\includegraphics[height=0.16\textheight]{obrazky/komp3.png}
}
\qquad
\subfloat[Pólenie cesty.]{
\label{img:komp4}
\includegraphics[height=0.20\textheight]{obrazky/komp4.png}
}
 
\label{img:kompresia}
\caption{\emph{Kompresia cesty z vrcholu 11 do koreňa.} 
Cesta je vyznačená šedou. 
Pri úplnej kompresii~\subref{img:komp2} sa všetky vrcholy 
napoja na zástupcu. Pri delení cesty~\subref{img:komp3} a pólení 
cesty~\subref{img:komp4} sa cesta skráti približne na polovicu.}
\end{figure}
 
Druhou heuristikou je kompresia cesty. Algoritmov na efektívnu kompresiu 
cesty je veľa \citep{paths2}, tu popíšeme tie najefektívnejšie. Prvou z nich 
je \emph{úplná kompresia} \citep{comp1}. Pri vykonávaní operácie $\find(x)$, po 
tom, ako nájdeme zástupcu, napojíme všetky vrcholy po ceste priamo pod koreň 
(obr.~\ref{img:komp2}). Toto síce trochu spomalí prvé hľadanie, ale výrazne 
zrýchli ďalšie hľadania pre všetky prvky na ceste ku koreňu. Druhou 
heuristikou je \emph{delenie cesty} \citep{comp2}. Pri vykonávaní operácie 
$\find(x)$ pripojíme každý vrchol v ceste od vrcholu $x$ po koreň stromu na otca 
jeho otca (obr.~\ref{img:komp3}). Treťou heuristikou je \emph{pólenie cesty} 
\citep{comp2}. Pri vykonávaní operácie $\find(x)$  pripojíme každý druhý vrchol 
v ceste od vrcholu $x$ po koreň stromu na otca jeho otca 
(obr.~\ref{img:komp4}).

Časová zložitosť union-findu záleží od toho, koľko prvkov je v množinách a koľko je 
operácií celkovo vykonaných operácií. Všetky uvedené spôsoby ako vykonať 
operáciu $\find(x)$ sa dajú použiť s obomi realizáciami operácie $\union$. 
Počet prvkov označme $n$ a počet operácií $m$. V praxi je zvyčajne počet 
operácií oveľa väčší ako počet prvkov. Pri tomto predpoklade ($m\geq n$) je 
pri použití spájania podľa ranku časová zložitosť pre algoritmus bez kompresie 
$\Theta(m\log n)$ a pre všetky tri uvedené typy kompresií 
$\Theta(m\mathop{\alpha}(m,n))$. 
V tabuľke~\ref{fig:uf:comp} je porovnanie časových zložitostí \citep{paths2}.

\begin{table}
\centering
\small
%\footnotesize %toto tu mám nechať?
\subfloat[Časová zložitosť pre union-find, keď $m \geq n$.][Prehľad časových zložitostí, ak $m \geq n$.]{
\begin{tabular}{m{5cm}m{4cm}m{4cm}}
& Naivné spájanie & Spájanie podľa ranku \tabularnewline
\hline
Naivné hľadanie & $\Theta\left( mn\right)$ & $\Theta\left( m\log n\right)$ \tabularnewline
Úplná kompresia, delenie cesty, pólenie cesty & $\Theta\left(m\log _{1+m/n} n \right)$ & $\Theta\left( m\alpha\left( m, n\right) \right)$ \tabularnewline
\label{fig:uf:comp1}
\end{tabular}
}
\qquad
\subfloat[Časová zložitosť pre union-find, keď $m < n$.][Prehľad časových zložitostí, ak $m < n$.]{
\begin{tabular}{m{5cm}m{4cm}m{4cm}}
& Naivné spájanie & Spájanie podľa ranku \tabularnewline
\hline
Naivné hľadanie & $\Theta\left( mn\right)$ & $\Theta\left(n + m\log n\right)$ \tabularnewline
Úplná kompresia& $\Theta\left(n + m\log  n \right)$ & $\Theta\left(n + m\alpha\left( n, n\right) \right)$ \tabularnewline
Delenie cesty & $\Theta\left(n \log m \right)$ & $\Theta\left(n + m\alpha\left( n, n\right) \right)$ \tabularnewline
Pólenie cesty & $\Omega\left(n + m\log n \right)$, $O\left(n\log m\right)$ & $\Theta\left(n + m\alpha\left( n, n\right) \right)$ \tabularnewline
\label{fig:uf:comp2}
\end{tabular}}
\caption{\normalsize Porovnanie časových zložitostí pre rôzne kombinácie hľadaní prvkov a 
spájaní množín pre union-find. Počet prvkov je $n$ a počet operácií je $m$. 
Ackermannova inverzná funkcia je označená ako $\alpha$. V praxi zväčša platí, že $m > n$.}
\label{fig:uf:comp}
\end{table}

% \section{Vizualizácia}
% Dátovú štruktúru union-find sme vizualizovali ako les. Pre názorné 
% oddelenie množín sme si zvolili pravidlo, ktoré zakazovalo vykresliť vrchol 
% z iného stromu (prvok inej množiny) 
% napravo od najľavejšieho vrcholu a naľavo od napravejšieho vrcholu inej 
% množiny. Jednotlivé množiny sme už vykreslovali tesným Walkerovým algoritmom 
% \citep{walker}. Vizualizácia poskytuje všetky vyššie spomínané heuristiky a 
% aj tlačidlo na vykonanie viacerých náhodných spojení naraz. Toto je užitočné, 
% keď chce uživateľ vidieť, ako dátová štruktúra vyzerá, po vykonaní 
% veľa operácií.
% 




\chapter{Písmenkový strom}\label{chap:trie}

\emph{Písmenkový strom} reprezentuje množinu slov. Oproti binárnym 
vyhľadávacím stromom je hlavný rozdiel v tom, že kľúče nie sú uložené 
vo vrcholoch, ale samotná poloha v strome určuje kľúč (slovo). 

\section{Popis}\label{sec:trie:popis}
Písmenkový strom je \emph{asociatívne pole (slovník)}, čiže 
poskytuje tieto tri operácie:
\begin{itemize}
\item $\put(\k)$ -- pridá do stromu slovo $\k$;
\item $\find(\k)$ -- zistí, či sa v strome slovo $\k$ nachádza;
\item $\delete(\k)$ -- odstráni zo stromu slovo $\k$.
\end{itemize}
Písmenkový strom je \emph{zakorenený strom}, v~ktorom každá hrana obsahuje 
práve jeden znak z abecedy alebo \emph{ukončovací znak}. Teda, každá cesta 
z koreňa do listu so znakmi $w_1, w_2, \ldots, w_n, \uz$ prirodzene 
zodpovedá slovu $w=w_1w_2\cdots w_n$. \emph{Ukončovací znak} je ľubovoľný, 
dopredu dohodnutý symbol, ktorý sa v abecede nenachádza (napr. \uz). 

Všetky tri operácie začínajú v koreni a ku slovu pridávajú ukončovací znak, 
teda pracujú s reťazcom $\k\uz$. 

\paragraph{Vkladanie slova.}%
Operácia $\put(\k)$ vloží do stromu vstupný reťazec tak, že z reťazca číta znaky 
a prechádza po príslušných hranách. Ak hrana s daným symbolom neexistuje, 
pridá ju (pozri obr.~\ref{img:trieinsert}).

\begin{figure}
\centering
\includegraphics[width=0.9\columnwidth]{obrazky/trieinsertsmall.png}
\caption{\emph{Vloženie slova \uv{{\tt SURNE}}.} Začiatok slova 
\uv{{\tt SU}} sa v strome nachádza, ešte treba pripojiť hrany 
so znakmi {\tt R}, {\tt N}, {\tt E} a \uz.} 
\label{img:trieinsert} 
\end{figure}

\paragraph{Hľadanie slova.}%
Operácia $\find(\k)$ sa spustí z koreňa podľa postupnosti znakov. Ak hrana, 
po ktorej sa má spustiť neexistuje, dané slovo sa v strome nenachádza. 
Ak prečíta celý vstupný reťazec, dané slovo sa v strome nachádza.

\paragraph{Mazanie slova.}%
Operácia $\delete(\k)$ najprv pomocou operácie $\find(\k)$ zistí umiestnenie slova. 
Ak sa slovo v strome nachádza, algoritmus odstráni hranu s ukončovacím 
znakom a vrchol, ktorý bol na nej zavesený. V tomto štádiu sa nám môže 
stať, že v strome ostane \emph{mŕtva vetva} -- nie je ukončená 
ukončovacím znakom. Pre fungovanie stromu to nevadí, všetky operácie by 
prebiehali správne, ale takto štruktúra zaberá zbytočne veľa miesta. 
Preto je dobré túto mŕtvu vetvu odstrániť (pozri obr.~\ref{img:triedelete}).

\begin{figure}
\centering
\includegraphics[width=0.9\columnwidth]{obrazky/triedeletesmall.png}
\caption{\emph{Odstránenie slova \uv{{\tt PODLA}}.} Po odstránení 
\uz\ nám v strome ostane nepotrebá prípona \uv{{\tt DLA}} (mŕtva vetva), 
ktorá je vyznačená červenou.}
\label{img:triedelete} 
\end{figure}

Všetky tri operácie majú časovú 
zložitosť $O(|\k|)$, kde $|\k|$ je dĺžka slova.

\section{Použitie}\label{sec:trie:pouzitie}
Prvýkrát navrhol písmenkový strom \citet{fredkin}, ktorý používal názov 
\emph{trie memory}\footnote{Z anglického re\emph{trie}val -- získanie.}, 
keďže išlo o spôsob udržiavania dát v pamäti. 
\citet{patricia} navrhol písmenkový strom, v ktorom sa každá cesta bez vetvení
skomprimuje do jedinej hrany (na hranách potom nie sú znaky, ale slová).
Táto štruktúra je známa pod menom \emph{PATRICIA} (tiež \emph{radix tree},
resp.\ \emph{radix trie}) a využíva sa napríklad v \emph{routovacích tabuľkách}
\citep{radix}.

Písmenkový strom (tzv.\ \emph{packed trie} alebo \emph{hash trie}) sa používa
napríklad v programe \TeX\ na slabikovanie slov \citep{liang}.
Pôvodný návrh \citep{fredkin} ako uložiť \trie\ do pamäte zaberal
príliš veľa nevyužitého priestoru. \citet{liang} však navrhol, ako
tieto nároky zmenšiť.

Písmenkové stromy sa podobajú na \emph{konečné automaty}. 
Vznikli rôzne modifikácie stromov na automaty, ktorých hlavnou výhodou je, 
že v komprimovanej podobe spájajú nielen predpony, ale aj prípony slov 
a teda v slovách prirodzených jazykov výrazne znižujú pamäťový priestor potrebný 
na uchovanie dátovej štruktúry. Vďaka tomu sa využívajú na jazykovú korekciu, 
automatické dopĺňanie slov a podobne \citep{scrabble,ca}. 

Ďalšie použitie písmenkového stromu je pri triedení zoznamu slov. 
Všetky slová sa pridajú do stromu a potom sa spraví \emph{preorderový prechod} 
stromu. Túto myšlienku spracovali \citet{burstsort1} a veľmi výrazne zrýchlili 
triedenie dlhých zoznamov slov. Neskôr tento algoritmus vylepšili 
\citet{burstsort2}. Algoritmus používa \trie\ s poľami, ktoré sa môžu zväčšiť. 
Pokiaľ sa pole zväčší za povolenú hranicu, "prepáli" sa. Preto bol 
algoritmus nazvaný \emph{burstsort}.

% Špeciálnym použitím písmenkového stromu je vytvorenie stromu zo všetkých 
% prípon slova. Táto dátova štruktúra sa nazýva \emph{sufixový strom} a dá sa 
% mo\-di\-fi\-ko\-vať na udržiavanie viacerých slov. Tieto štruktúry majú 
% veľmi veľa praktických využití \citep{gusfield}. 

% \section{Vizualizácia} 
% Pri vizualizácií písmenkového stromu sme použili 
% Walkerov algoritumus pre úsporné rozloženie vrcholov v strome
% \citep{walker}. Keď má vrchol viacej synov a hrany kreslíme priamo, tak vzniká 
% nedostatok priestoru pre umiestnenie znakov na hrany. Preto sme sa rozhodli 
% kresliť hrany zakrivene, podľa Bézierovej krivky určenej štyrmi bodmi. 
% Vo vizualizácií sa dajú vložiť náhodné slová podľa momentálne nastaveného 
% jazyka. Taktiež sa automaticky odstraňuje diakritika a interpunkcia, takže 
% sa dá naraz vložiť súvislý text.



\begin{figure}[!hbt]
\begin{center}
\includegraphics[width=0.5\textwidth]{obrazky/komlogo-new}

\caption{Ján Ámos Logo.}
\label{img:komensky}
\end{center}
\end{figure}

A na Obrázku \ref{img:komensky} je väčší Komenský ako na titulke.

\begin{verbatim}
for(int i=0; i<n; i++)
    cout << "loop";
\end{verbatim}

\begin{figure}[!hbt]
\begin{center}
\includegraphics[width=\textwidth]{obrazky/uf00.png}

\caption{Union-find procedúra find - hľadanie zástupcu.}
\label{img:uf00}
\end{center}
\end{figure}

\begin{align}
& \cos^2\psi+\sin^2\psi & = 1, \\
& \cosh^2\omega-\sinh^2\omega &=1. \\
& ⊢ ∀x[(Fx ∨ Gx) → \mathord{∼}Hx] \\
& ⊨ ¬∃y∀x[x∈y ↔ x∉x]  \\
& ⊭ x ∩ (y ∪ z) ≠ (x ∩ y) ∪ (x ∩ z) \\
& ⊢ ⟦α⟧ = ℵ₀ → α ≇ ℘(α) \\
& ⌜ψ[(℩x)φx]⌝ ≝
     ⌜(∃x)[φx ∧ (∀z)(φz ⊃ x=z) ∧ ψx]⌝ \\
& ⊢ (P ⥽ Q) ⥽ (□P ⥽ ◇Q)
\end{align}

{\fontspec[LigType=disc]{Junicode}[20][[1]]}


\cleardoublepage
\phantomsection
\addcontentsline{toc}{chapter}{Záver}
\chapter*{Záver}\label{chap:zaver}

V práci sme popísali dátové štruktúry union-find, písmenkový strom a sufixový 
strom (kapitoly~\ref{chap:uf}, \ref{chap:trie} a \ref{chap:sx}). Spísali sme 
návrh a implementáciu softvéru. Softvér sme implementovali 
podľa návrhu a splnili sme všetku funkcionalitu: vizualizovali sme všetky 
dátové štruktúry a algoritmy, doplnili funkčnosť, spravili komentáre v 
angličtine a slovenčine.

V ďalšom vývoji by sme chceli prerobiť komentáre pre sufixové stromy, keďže sa 
nám zdajú byť neprehľadné. Taktiež by sme chceli zvýrazniť niektoré časti 
Ukkonenovho algoritmu a pridať nové algoritmy vymenované v sekcií 
\ref{sec:sx:usage}.

Softvér by sme chceli v budúcnosti rozšíriť o nové dátové štruktúry, ktoré sa 
týkajú stringológie, napríklad \emph{radix tree}, \emph{PATRICIA}, 
\emph{orientovaný acyklický graf pre slová (DAWG)}.

Chceli by sme projekt podstúpiť širšiemu použitiu a zozbierať viac spätnej 
väzby, aby sme softvér odladili podľa používateľských potrieb a opravili 
chyby, ktoré sme neobjavili. 


\cleardoublepage
\phantomsection
\addcontentsline{toc}{chapter}{Literatúra}

\printbibliography

\backmatter

\end{document}
